\introduction % Структурный элемент: ВВЕДЕНИЕ

Актуальность темы данной работы связана с растущей нагрузкой на детских психологов в образовательных учреждениях и медицинских центрах. Рисуночные тесты --- одна из наиболее распространённых проективных методик, применяемых для предварительного скрининга психологических состояний детей дошкольного и младшего школьного возраста. Среди них особое место занимает тест <<Рисунок человека>>, методология интерпретации которого детально разработана А.~Л.~Венгером. Однако ручная обработка рисунка специалистом занимает значительное время и подвержена субъективным факторам. Автоматизация хотя бы части этого процесса с помощью методов машинного обучения позволяет масштабировать применение методики и обеспечить воспроизводимость результатов.

Современные нейронные сети успешно применяются для широкого спектра задач классификации изображений. Однако специфика детских рисунков --- бледные карандашные штрихи, схематичность, разнообразие стилей, разная плотность заполнения листа --- ставит ряд особых задач: точная локализация фигуры на странице, обработка изображений с низким контрастом, корректное распознавание сильно несбалансированных категорий признаков.

Помимо технической задачи распознавания признаков, существует не менее важная задача доменной интерпретации: нейросеть выдаёт техническую характеристику рисунка, например: <<размер крупный>>, <<нажим сильный>>, <<отсутствует шея>> --- но конечному пользователю необходимо получить психологически осмысленный текст, обобщающий совокупность признаков в характеристику ребёнка.

Цель работы --- разработать нейросетевую систему распознавания визуальных признаков детских рисунков и их сопоставления с психологическими чертами на основе методики А.~Л.~Венгера. Для достижения поставленной цели в работе были решены следующие задачи:

\begin{itemize}
    \item проведён анализ предметной области рисуночных тестов и существующих подходов к автоматизации их интерпретации;
    \item разработан алгоритм предобработки изображений, реализующий локализацию фигуры человека на листе на основе классических операций обработки изображений;
    \item спроектирована и обучена собственная нейронная сеть для одновременной классификации признаков рисунка;
    \item составлен словарь интерпретаций, сопоставляющий технические признаки рисунка с каноническими психологическими чертами по А.~Л.~Венгеру;
    \item реализован модуль генерации текстового психологического портрета на основе предсказаний модели;
    \item проведено экспериментальное исследование качества модели на тестовой выборке;
    \item для удобства практического использования разработанные модули объединены в микросервисный веб-сервис, обеспечивающий доступ к функциональности системы через браузер.
\end{itemize}

Для разработки системы необходимо изучить инструменты и методы, решающие поставленные задачи. Работа основывается на следующих библиотеках, технологиях и алгоритмах:

\begin{itemize}
    \item Python является основным языком программирования, применяемым при разработке;
    \item PyTorch обеспечивает построение, обучение и применение нейронных сетей;
    \item OpenCV применяется для реализации алгоритма локализации фигуры на изображении;
    \item scikit-learn обеспечивает разбиение датасета и расчёт метрик качества;
    \item iterative-stratification используется для группового стратифицированного разбиения мультиметочных данных;
    \item Pillow выполняет загрузку, преобразование и сохранение изображений;
    \item Google Colab Notebooks предоставляет вычислительную платформу с графическим ускорителем для обучения модели;
    \item FastAPI, PostgreSQL, MinIO и Docker используются для оформления решения в виде микросервисного веб-сервиса.
\end{itemize}

В результате выполнения работы был разработан программный комплекс на языке Python, состоящий из модуля препроцессинга, обученной нейросетевой модели, словаря интерпретаций признаков и генератора текстового психологического портрета. Дополнительно эти компоненты объединены в микросервисную веб-систему, развёртываемую через docker-compose, что обеспечивает доступ к функциональности через браузер без необходимости установки локального окружения.

Результаты работы могут быть использованы психологами и педагогами в образовательных и медицинских учреждениях для предварительного анализа, а также в научных исследованиях для автоматизированной обработки больших объемов детских рисунков.
